\section{Background}
\label{sec:background}

The emergence of large language model--driven agents has shifted the unit of computation from isolated inference calls to autonomous processes that perceive, plan, act, and spawn sub-agents across distributed environments. As these multi-agent systems (MAS) grow in autonomy and scale, foundational concerns previously studied in distributed systems---lifecycle management, accountability, provenance, and hierarchical decomposition---resurface with new urgency. This section situates recursive spawning within that lineage, reviewing agent lifecycle management in MAS (\S\ref{sec:lifecycle}), accountability and provenance models for distributed agents (\S\ref{sec:provenance}), the distinction between fixed and mutable delegation chains (\S\ref{sec:delegation}), and hierarchical task decomposition as a substrate for immutable parent chains in the BX3 framework (\S\ref{sec:bx3_substrate}).

\subsection{Agent Lifecycle Management in Multi-Agent Systems}
\label{sec:lifecycle}

Classical distributed systems treat an agent's lifecycle as a managed state machine: creation, activation, execution, suspension, migration, and termination are first-class events governed by explicit protocols~\cite{ferrarotti2022lifecycle}. These protocols ensure that system resources are allocated coherently, that agents can be tracked across distributed nodes, and that cleanup procedures fire reliably on exit. The state-machine model presupposes that agent behavior is deterministic given its program and inputs---an assumption that holds for traditional software agents but breaks down for LLM-based agents whose behavior includes probabilistic reasoning and adaptive planning.

Modern LLM-based MAS extend the classical lifecycle model in two non-trivial directions. First, agents are no longer statically programmed but dynamically instantiated from prompts, tool manifests, and contextual embeddings, blurring the line between code and execution context. A spawned sub-agent inherits not only a task description but a partial reasoning state drawn from its parent's conversational context, creating a heritage of assumptions that may not be explicitly represented in any formal contract between parent and child. Second, agents can now spawn sub-agents that themselves spawn further agents, forming emergent topologies---trees, forests, or star graphs---that are not known at design time and may exceed the intentional depth of the original system architecture~\cite{agenthive2025}.

Several frameworks have attempted to address this extended lifecycle model. AgentOrchestra introduces a hierarchical multi-agent framework in which a central planning agent decomposes complex objectives into sub-goals and delegates them to specialized sub-agents, each equipped with domain-specific tooling~\cite{agentorchestra2025}. The framework implements lifecycle management through three coordinated stages: a planning phase where the orchestrator analyzes the task graph, a delegation phase where sub-agents are instantiated with scoped permissions, and a synthesis phase where results are aggregated and fed back to the planner. The planning agent plays the role of a lifecycle controller---it monitors sub-agent health, re-decomposes on failure, and retires agents whose tasks are complete. This controller pattern mirrors the Manager/Worker architecture in distributed computing, but with the added complexity that LLM-based agents are probabilistic and can deviate mid-execution, requiring built-in verification gates that classical lifecycle models do not specify.

AGENTHIVE generalizes the controller pattern by making delegation itself a first-class primitive rather than a side effect of orchestration~\cite{agenthive2025}. In AGENTHIVE, any agent can spawn, coordinate, and retire sub-agents via first-class delegation operations, enabling decentralized control without a central orchestrator. Recursive delegation forms task-driven emergent topologies---trees, forests, and star formations---without any pre-defined agent graph. Lifecycle stages emerge dynamically: agents autonomously create sub-agents as tasks demand, monitor their progress through explicit signaling, and dissolve them upon completion or failure. This model aligns lifecycle management with the task structure itself rather than with a predetermined organizational chart, which is better suited to open-ended workloads where the optimal topology is discovered rather than prescribed.

A complementary line of work has focused on formalizing the lifecycle stages themselves. The Unified Agent Lifecycle Management (UALM) model proposes five layers of lifecycle governance---provisioning, authorization, execution, monitoring, and retirement---with explicit state transitions defined at each layer~\cite{ualm2026}. The Agent Lifecycle Protocol (ALP) working group has begun standardizing these states and transitions as a formal specification for multi-agent systems interoperability, drawing on analogies to transport-layer protocol design where well-defined state machines prevent resource leaks and ambiguous authority roles~\cite{alp2025}. Together, these efforts establish that lifecycle management in LLM-based MAS is not merely an engineering concern but a formal specification problem: without explicit, machine-readable lifecycle semantics, sub-agents can execute beyond their authorized window without any systemic means to detect or terminate them. This is not a hypothetical edge case; it is the observed behavior of early multi-agent deployments where long-running agent loops accumulated orphaned child processes that continued consuming compute while producing no aligned output.

The broader implication is that lifecycle management for autonomous agents must address two concurrent requirements: the structural requirement to maintain a complete, verifiable ancestry graph from any active agent back to an authorizing principal, and the operational requirement to allow that graph to evolve dynamically as tasks decompose and recombine. These requirements are in tension with each other, and resolving that tension is the central design challenge for any spawning architecture.

\subsection{Accountability and Provenance in Distributed Agent Systems}
\label{sec:provenance}

A spawned sub-agent acts on behalf of its parent, which itself may be acting on behalf of another agent or a human. When something goes wrong---when an agent produces a harmful output, makes an unauthorized decision, or passes corrupted state downstream---the first questions any audit asks are: who authorized this action? Which agent in the chain introduced the error? How does one construct a forensically reliable audit trail when agents are probabilistic and may not retain complete records of their own reasoning? These questions are not merely technical; they are questions of accountability architecture.

The Human Delegation Provenance (HDP) protocol addresses this by introducing a lightweight, token-based system that records and cryptographically proves the chain of authorization from a human principal through all downstream agents~\cite{hdp2026}. Each hop in the delegation chain is recorded as a signed attestation, building an append-only chain that any verifier can check offline using only the issuer's public key. Crucially, HDP separates the identity of the human authorizer from the execution chain, enabling post-hoc reconstruction of intent even in deeply nested agent graphs. The protocol handles multi-hop delegation where agents delegate to sub-agents, which in turn invoke tool-execution agents, preserving the chain of custody across all hops. The design choice to make the chain append-only is deliberate: it ensures that even a compromised or malfunctioning agent cannot erase the record of what it was authorized to do.

The Warrant Certificate Authority (WCA) framework approaches the accountability problem from the data-integrity side, focusing on the provenance of information as it flows through tool-call chains~\cite{sentinelagent2026}. WCA introduces attestation structures that certify the origin and trust pathway of data at each tool-call boundary, directly addressing what the framework calls ``semantic laundering''---the phenomenon where data inherits unwarranted credibility as it crosses tool boundaries. A PKI-like chain of certificate authorities establishes provenance assurances at each hop, and Warrant Attestation Levels (WAL-0 through WAL-3) provide incremental adoption paths analogous to SLSA build levels in software supply chains~\cite{sentinelagent2026}. The WAL framework is particularly relevant for recursive spawning because it demonstrates that trust verification can be layered: systems can begin with lightweight attestation and increase rigor as stakes grow, without requiring a complete redesign of the trust architecture.

CollectiveOS v4 operationalizes these concerns in a production multi-agent operating system, introducing a cryptographic ledger called the Proof Chain that binds agent intent, computational lineage, and governance policy into immutable receipts for every micro-transaction~\cite{collectiveos2025}. The system uses OpenTelemetry-based observability to capture system-wide telemetry, enabling end-to-end auditing of agent reasoning, state changes, and governance checks at every layer of the stack. A Forensic Proof Vault stores evidence of decisions and actions, strengthening chain-of-custody claims for regulatory and legal proceedings. This represents the most complete operationalization of provenance as a first-class construct in agentic systems to date: the system does not merely allow auditing, it structures every agent interaction as a potential evidence artifact.

PROV-AGENT extends the W3C PROV data model with agent-centric metadata---prompts, responses, decisions, and intermediate reasoning states---linked to broader workflow context and downstream outcomes~\cite{proveagent2025}. The framework captures provenance across federated, heterogeneous environments spanning edge devices, cloud infrastructure, and high-performance computing clusters, and provides tooling for near real-time provenance queries. The motivating problem is the cascading error problem: a hallucination or reasoning error in one agent can propagate, undetected, through subsequent agents in the chain, compounding into a system-level failure whose origin is difficult to isolate. PROV-AGENT addresses this by making each agent's provenance record a first-class, queryable object, enabling auditors to trace the propagation of any error back to its origin point.

These frameworks collectively establish that accountability in MAS requires three structural properties that cannot be satisfied by policy or convention alone. First, \emph{traceability}: every action must be traceable to an authorizing principal, and that trace must be machine-verifiable rather than relying on self-reported logs. Second, \emph{tamper-evidence}: the audit trail must be resistant to retrospective modification by any actor in the system, including the agent whose behavior is being audited. Third, \emph{compositional verification}: it must be possible to verify the full delegation chain without relying on centralized infrastructure, since centralized infrastructure is itself a single point of failure and a scalability bottleneck.

\subsection{Fixed versus Mutable Delegation Chains}
\label{sec:delegation}

A foundational design choice in any spawning system is whether the delegation chain is fixed at creation or mutable during execution. A \emph{fixed delegation chain} establishes a static parent--child relationship at spawn time: the sub-agent inherits a fixed identity, a fixed scope, and a fixed communication channel back to its parent. This model simplifies accountability because each hop in the chain is known in advance and cannot change; forensic reconstruction reduces to traversing the pre-established parent links. The tradeoff is rigidity: a fixed chain cannot adapt to task evolution, agent failure, or scope drift without explicit re-authorization from the parent.

Mutable delegation chains allow the chain of authority to be reconfigured at runtime. Agents may re-delegate tasks, re-scope permissions, or spawn additional sub-agents not anticipated at the time of the original request. Mutable chains support adaptive task redistribution and fault tolerance---if a sub-agent fails, its task can be reassigned to a peer without waiting for the original parent to intervene. However, mutable chains introduce forensic complexity: when a chain is modified mid-execution, the original authorization context may be lost or blurred. An auditor examining a mutable chain after a failure must reconstruct not only who authorized what, but which version of the chain was in effect at the time of the adverse event.

The HDP protocol's append-only hop record provides a principled resolution to this tension: even mutable chains can maintain an immutable provenance log of every change, preserving the original authorization context alongside runtime modifications~\cite{hdp2026}. Under this model, the chain's topology may change during execution, but its complete history is preserved. Each modification is recorded as a new hop with a timestamp, a signer, and a scope modification record. This turns the forensic challenge from reconstructing a lost authorization state into traversing an append-only log---a structurally simpler problem.

The SentinelAgent framework formalizes the requirements for delegation chain integrity through the Delegation Chain Calculus (DCC), which defines seven properties that any delegation chain---fixed or mutable---must satisfy to be considered secure and accountable~\cite{sentinelagent2026}. Among these, \emph{authority narrowing} ensures that delegated authority never exceeds the grantor's authority at the time of delegation; \emph{forensic reconstructibility} guarantees that any action can be traced to the human who authorized the original intent; and \emph{scope-action conformance} ensures that each agent's actions remain within the scope declared at the time of delegation. The calculus demonstrates that mutable chains can satisfy these properties if and only if every modification is recorded as a new hop in the delegation chain, never by editing or deleting prior hops. This constraint effectively reconciles mutability with auditability: the chain's topology may change, but its history cannot. A fixed delegation chain can be viewed as the constrained special case of this model where the mutation log is empty from the outset. The advantage of the fixed model is simplicity and performance---no hop-resolution overhead during runtime traversal. The advantage of the mutable model is flexibility and fault tolerance. The BX3 framework, described next, resolves this tension by treating the delegation chain as immutable at the level of parent identity while allowing the child agent full autonomy over its own execution and further spawning.

\subsection{Hierarchical Task Decomposition and Immutable Parent Chains in the BX3 Framework}
\label{sec:bx3_substrate}

Hierarchical task decomposition (HTD) is the mechanism by which a complex goal is recursively partitioned into sub-goals assigned to specialized agents. HTN planning provides the classical formal foundation: a complex task is decomposed into simpler subtasks via explicitly defined task-reduction rules, and the decomposition continues until subtasks are primitive enough to execute directly~\cite{ghallab2004htn}. This body of work establishes that hierarchical decomposition is not merely a convenient engineering pattern but is, in the bounded rationality tradition, a necessary condition for tractable planning in real-world environments where the state space is too large for flat search to be feasible~\cite{erpl,simon}.

HALO introduces a three-layer hierarchy for multi-agent LLM systems in which a high-level planner decomposes tasks, mid-level role-design agents instantiate sub-agents for each subtask with appropriate tool access and constraints, and low-level inference agents execute concrete steps~\cite{halo2025}. The decomposition is not fixed: the planner can dynamically re-instantiate agents as the task evolves, using Monte Carlo Tree Search to explore optimal reasoning trajectories. This adaptive decomposition is critical in open-ended tasks where the optimal sub-task graph is not known upfront and must be discovered through interleaved planning and execution.

ReAcTree extends this to long-horizon task planning by arranging LLM-based agents in a dynamic tree structure with explicit control-flow nodes that coordinate strategy across agent nodes~\cite{reactree2025}. Each agent node can independently reason, act, and further expand the tree---enabling recursive decomposition where a sub-agent's own planning process spawns additional sub-agents. The tree expansion is gated by a convergence criterion that prevents unbounded growth: once a subtree's expected remaining work falls below a threshold, further spawning stops and the results propagate upward. Memory is split into episodic (goal-specific examples) and working (environment-specific observations) components, enabling each node to remain aligned with the global objective while maintaining local autonomy. This separation of episodic and working memory is architecturally important for recursive spawning because it prevents unbounded context growth while preserving the reasoning trace necessary for accountability.

The BX3 framework adopts hierarchical task decomposition as its foundational execution model but introduces a structural constraint that differentiates it from all prior work in this space: \emph{immutable parent chains}. In BX3, every spawned agent records the identity of its parent at spawn time, and this parent identity is cryptographically immutable---neither the child agent nor any subsequent actor in the system can modify or obscure the parent link after creation~\cite{bx3-framework}. This immutable parent chain serves as a forensic backbone: any action taken by any agent in the system can be traced, through recursively following parent links, back to a specific root agent or human principal. The chain is not stored as a reference in the agent's memory (which is mutable and unreliable) but is embedded in the agent's identity record at the kernel level of the AgentOS runtime.

This design draws on the provenance guarantees established in HDP and WCA but applies them at the structural level rather than as an overlay. Rather than relying on a separate attestation log that must be actively maintained and reconciled against runtime events, the parent chain itself provides the audit trail: each agent's parent identity is embedded in its execution context and verifiable at runtime without centralized lookup. The immutable chain also enables a property BX3 calls \emph{scope inheritance with attenuation}: a child agent inherits the operational constraints of its parent but may narrow them; it may not widen them. This ensures that authority never escalates through delegation---a property that SentinelAgent's DCC formally proves is necessary for secure multi-agent execution~\cite{sentinelagent2026}. The combination of hierarchical decomposition (adapted from HALO and ReAcTree's planning model) with immutable parent chains (BX3's structural constraint) provides a substrate for recursive spawning that is simultaneously flexible---agents can spawn as needed to match task structure---and accountable---every spawn event is anchored to an immutable parent, and every action can be traced to its origin without relying on self-reported logs.

